\documentclass[11pt,letterpaper]{article}
\usepackage[T1]{fontenc}
\usepackage[utf8]{inputenc}
\usepackage{lmodern,geometry,microtype,xcolor,enumitem,hyperref,fancyhdr}
\geometry{margin=0.78in}
\definecolor{ink}{HTML}{17324D}
\definecolor{accent}{HTML}{126A78}
\hypersetup{colorlinks=true,linkcolor=accent,urlcolor=accent,breaklinks=true}
\setlength{\parindent}{0pt}
\setlength{\emergencystretch}{3em}
\setlist[itemize]{leftmargin=*,itemsep=3pt,topsep=3pt}
\pagestyle{fancy}\fancyhf{}\fancyhead[L]{\small Architecture Documentation Suite}\fancyhead[R]{\small v1.0}\fancyfoot[R]{\thepage}\setlength{\headheight}{15pt}
\begin{document}
\thispagestyle{empty}
{\color{accent}\large TECHNICAL REFERENCE \hfill NOVEMBER 2026}\par\vspace{1.1cm}
{\color{ink}\LARGE\bfseries\raggedright Views and Beyond: Technical Study}\par\vspace{0.35cm}
{\Large Architecture structures, documentation economics, and engineering practice}\par\vspace{0.55cm}
\rule{\textwidth}{0.7pt}\par\smallskip
{\small Architecture Documentation Suite | Release 1.0}\par\vspace{0.7cm}
\tableofcontents
\newpage
\section*{1. Purpose and source boundary}\addcontentsline{toc}{section}{1. Purpose and source boundary}
The Views and Beyond (V\&B) approach makes architecture documentation useful for implementing, maintaining, communicating and analyzing software. Its central organizing decision is to describe relevant views and then record information that applies across those views. This study synthesizes Clements et al., second edition, rather than reproducing the book. Source: Prologue P.1--P.6, chapters 1--11, Epilogue and Appendix A. The source uses earlier architecture-description standard references; current-edition mapping belongs in document 04.\par\smallskip
Architecture is not the diagram or document: it is the set of significant structures of the system. A structure consists of elements and relationships; a view is its representation for a purpose. A style supplies an architectural vocabulary, including permitted elements, relationships, properties and constraints. A pattern adds problem context and recurring solution guidance. A system may exhibit several styles simultaneously.\par\smallskip
\section*{2. Documentation principles and economics}\addcontentsline{toc}{section}{2. Documentation principles and economics}
\begin{itemize}
\item Reader orientation: identify specific stakeholder work and write to enable it; a developer and an operator may need different views.
\item Minimize duplication: keep a canonical catalog and reference stable identifiers in other views.
\item Eliminate ambiguity: define the meaning, direction and constraints of every relationship; do not assume arrows mean calls.
\item Standardize organization: use the same view packet layout and interface record across projects.
\item Preserve decision rationale: distinguish alternatives, assumptions, tradeoffs and rejected options from the observed implementation.
\item Maintain useful currency: version against releases or commits; mark stale content rather than pretending it is current.
\item Review for fitness: ask representative stakeholders to perform tasks using the description, not merely approve the aesthetics.
\end{itemize}
Documentation is an investment: choose detail in proportion to maintenance cost and decision value. Record significant changes as they happen and publish at useful baselines. A view with no audience or analytical purpose is a candidate for omission. [Source P.2, P.5, E.4.]\par\smallskip
\section*{3. Module structures: implementation units}\addcontentsline{toc}{section}{3. Module structures: implementation units}
Module views answer how software responsibilities and implementation dependencies are organized. A module is not necessarily a running process. Document module responsibility, interfaces, allowed dependencies, visibility and relevant implementation properties.\par\smallskip
\begin{itemize}
\item Decomposition: nested responsibilities, ownership and part-of semantics; useful for change impact and work partitioning.
\item Uses: correctness depends on another unit's correct operation; do not equate every import with a demonstrated uses relationship.
\item Generalization: specialization and substitutability of abstractions; show contracts and invariants.
\item Layered: permitted dependency directions; state whether skipping layers or same-layer calls are allowed.
\item Aspects: crosscutting concern mechanisms and where they apply.
\item Data model: entities, cardinalities, constraints, ownership and lifecycle; distinguish logical data model from storage layout.
\end{itemize}
A dependency structure matrix can reveal cycles. For a layered system, construct an allowed-dependency matrix and compare it with static evidence. For a hexagonal design, document ports in the core and adapter dependencies separately from runtime invocation direction. [Source chapters 1--2.]\par\smallskip
\section*{4. Component-and-connector structures: runtime interaction}\addcontentsline{toc}{section}{4. Component-and-connector structures: runtime interaction}
C\&C views describe executing units and their interactions. Record components, ports, connectors, roles, protocols and relevant connector properties. Distinguish component types from instances; a single module can participate in several runtime components.\par\smallskip
\begin{itemize}
\item Pipe-and-filter: transformations and flow of data; document formats, buffering and backpressure.
\item Client-server: requests and responses; document endpoints, availability, timeouts and failure semantics.
\item Peer-to-peer: participants with peer roles and coordination rules; document discovery and failure recovery.
\item Service-oriented: independently exposed services and contracts; capture invocation, ownership and policy boundaries.
\item Publish-subscribe: publishers, event channels and subscribers; record delivery, ordering and idempotency semantics.
\item Shared-data: cooperating components using shared storage; record consistency, concurrency and data ownership.
\end{itemize}
Tiers concern runtime organization and should not be conflated with module layers. A runtime view can represent dynamic component creation if the cardinality, creation trigger and termination rule are stated. [Source chapters 3--4.]\par\smallskip
\section*{5. Allocation structures: software to environment}\addcontentsline{toc}{section}{5. Allocation structures: software to environment}
Allocation views connect software to nonsoftware structures. Deployment maps runtime elements onto execution nodes and networks; install maps software artifacts to installation locations; work assignment maps implementation units to responsible organizational groups. Other useful allocations may include requirements, resource budgets and ownership. Do not infer a live deployment merely from a Dockerfile or Terraform configuration. Document intended topology separately from observed deployment. [Source chapter 5.]\par\smallskip
\section*{6. Beyond the structures}\addcontentsline{toc}{section}{6. Beyond the structures}
A context diagram places the subject within its environment using the vocabulary of the view. Refinement can expose decomposition or additional design detail. Variation points record allowed alternatives, binding time, mechanism and constraints. A decision record should state the issue, options, criteria, resolution, status and consequences; only infer rationale if explicitly labeled as inference. Combine views when their concepts remain clear and a stakeholder benefits; otherwise cross-reference them. [Source chapter 6.]\par\smallskip
Interface documentation should communicate syntax and semantics: provided/required service, inputs, outputs, preconditions, postconditions, errors, sequencing, restrictions, quality properties and evolution policy. A signature alone is insufficient. Behavior can be specified as representative traces or more comprehensive models; choose based on the questions to answer. Use a sequence diagram for an interaction trace, a state machine for lifecycle rules and an activity diagram for branching workflows. [Source chapters 7--8, Appendix A.]\par\smallskip
\section*{7. View selection and packaging}\addcontentsline{toc}{section}{7. View selection and packaging}
Select candidates by mapping stakeholder concerns to structures; consolidate overlapping candidates when clarity and coverage survive; prioritize production according to stakeholder need and cost. A useful default for an existing service is context, module decomposition, runtime interaction and deployment, with behavior and security added when warranted. That default is a local starting point, not a universal prescription. [Source chapter 9.]\par\smallskip
Each view packet contains a primary presentation, catalog of elements and relations (including relevant interfaces and behavior), context, variability and rationale. Beyond views, provide document control, roadmap, documentation conventions, system overview, cross-view mapping, rationale and directory. Trace requirements where needed. [Source chapter 10.]\par\smallskip
\section*{8. Review and worked example}\addcontentsline{toc}{section}{8. Review and worked example}
For a job-processing service, developer concern C-01 is change impact, addressed by a module view showing API, domain and adapters. Operator concern C-02 is recovery, addressed by a runtime view and deployment view; the retry state machine documents failure behavior. Map the API module to the HTTP process, domain modules to the worker process and persistence adapter to the storage connector. These mappings are illustrative, not facts about a real repository.\par\smallskip
Run a review against a fixed baseline: establish purpose and subject, tailor question sets, plan participants and review logistics, conduct the review, then analyze results. Ask the operator to locate failure recovery, the developer to trace a proposed change and the security reviewer to identify authority boundaries. Record omissions and assign corrective owners. A successful document review does not prove the architecture itself is sound. [Source chapter 11.]\par\smallskip
\section*{9. Practical adoption}\addcontentsline{toc}{section}{9. Practical adoption}
Start with a thin baseline and add view packets when decisions require them. Keep stable element IDs, record the commit SHA, use text-based diagram sources and verify relationships against evidence. Gate changes when significant architecture facts drift. For the concrete implementation workflow see documents 02, 04, 05 and 06; for agent runtime specialization see 07.\par\smallskip
\section*{Bibliographic foundation}
Paul Clements et al., \emph{Documenting Software Architectures: Views and Beyond}, 2nd ed., Addison-Wesley/Pearson, copyright 2011, ISBN 978-0-321-55268-6. Section and chapter references refer to the supplied edition. Additional framework sources are identified in the applicable documents.\par
\end{document}
